\section{Actual formal degrees and a direct construction of the level-prime size bound}
\label{sec:degree-prime-size}

\status{DSZ1--DSZ3 close the original two-sided degree input $h_3$ on
every supplied family, using the already matched Harish--Chandra input
of DC18. DSZ4--DSZ9 prove the literal \texttt{TuSize} from the original
eventual volume identity. DSZ10 supplies that identity on the original
\texttt{Built} construction class and proves its \texttt{Closed} size
obligation. The common class and its actual families still have to be
constructed. The proof does not require the stronger Vandermonde
variation or prime-counting package.}

Write $W_n=\prod_{j=0}^{n-1}j!$ and
$V_n(z)=\prod_{a<b}(z_a-z_b)$. All ranks, coordinates, levels and
volumes below belong to the same original datum. In the family
calculation $L=x^\epsilon$, $B=e^L$, $a=xB$, $r=k+m-1$, and
$f=[E:\mathbb Q]/2$. The polynomial dimension is $d_n(z)=V_n(z)/W_n$.
The letter $\delta$ in DSZ2 is an error tolerance, not the family
exponent $\epsilon$ and not a Harish--Chandra coordinate gap.

% BEGIN DSZ1
\paragraph{DSZ1: the finite product comparison with the exact compact type.}
\label{claim:DSZ1}
Let $\lambda_0>\cdots>\lambda_{n-1}$ be regular integral in the sense
of (DC.1), and let $s$ be the original sign vector. Put
$p=\#s^{-1}(+)$, $q=\#s^{-1}(-)$, and
$A=\max_{a,b}|\lambda_a-\lambda_b|$, with $A=0$ for an empty
index set. Define
\[
 d_\tau(s,\lambda)=
 \frac{\displaystyle\prod_{a<b,\ s_a=s_b}
       (\lambda_a+\rho_{n,a}-\lambda_b-\rho_{n,b})}{W_pW_q}.
                                                               \tag{DSZ.1}
\]
Here $\rho_n$ is the half-sum of the noncompact positive roots,
with exactly the convention of DC1. If $a<b$ and $s_a=s_b$, direct
counting in its definition gives
\[
 \rho_{n,a}-\rho_{n,b}
       =\#\{c:a<c<b,\ s_c\ne s_a\}\ge0.
\]
Thus each compact numerator factor in (DSZ.1) dominates the
corresponding factor of $V_n(\lambda)$. There are exactly $pq$
unordered pairs of opposite signs; choosing their increasing order
identifies them with the noncompact positive roots. Each of their
factors is at most $A$. Finally
\[
 W_pW_q\le W_{p+q},
\]
because $W_{p+q}/W_p=\prod_{j=0}^{q-1}(p+j)!\ge\prod_{j=0}^{q-1}j!$.
Multiplying and dividing these positive quantities proves
\[
             0<d_n(\lambda)\le d_\tau(s,\lambda)A^{pq}.
                                                               \tag{DSZ.2}
\]
For $n=0$ both products are empty, $p q=0$ and the assertion is
$1\le1$; for $n=1$ the same convention applies. For $n\ge2$,
regular integrality makes every positive difference an integer,
so $A\ge1$. DC1 identifies (DSZ.1) with the dimension of the
specified compact type, and DC10 identifies the actual lowest
compact-type subspace with that type in the constructed large-gap
representation. DC13 transfers this identification to every literal
\texttt{IsDS} representation in the same large-gap range.
\emph{Source: Z01--Z03.}
% END DSZ1

% BEGIN DSZ2
\paragraph{DSZ2: the source gap reaches the reconstructed existence range.}
\label{claim:DSZ2}
Suppose $p,q\ge1$, $n=p+q$, $e^{-n}\le\delta<1/2$, and
\[
 h\ge8192\min(p,q)(pq)^3\frac{\log n}{\delta^2}.
                                                               \tag{DSZ.3}
\]
Then $\delta>0$, $n\ge2$, $pq\ge1$, $n\le2pq$, and
$\log n\ge\log2>1/2$. It follows that
$h>16384(pq)^3>32npq+2$: indeed
$32npq+2\le64(pq)^2+2\le66(pq)^3$.
Writing $q_0=\min(p,q)$, we have
$\sqrt{2q_0}\le2q_0$, so
\[
 16nq_0+8\sqrt{2q_0}\le24nq_0\le24npq<32npq+2.
\]
The other threshold $npq+1$ is smaller too. Consequently (DSZ.3)
implies the actual existence threshold (DC.2), and also the smaller
rigidity threshold used in DC13. This checks the needed hypothesis
against the reconstructed DC argument, rather than only against the
formal source's potentially smaller existence bound.

The general construction of DC1--DC12 therefore supplies a genuine
representation at this parameter, with its entire literal
\texttt{IsDS} predicate as proved at the end of DC12.
DC18--DC19 identify its degree, for the exact compact-dual Haar
measure, with $d_n(\lambda)$. This uses the already stated published
input HC-Deg, whose group-and-Haar constant DC16--DC18 determine.
The input applies to the connected finite-center group $SU(p,q)$;
DC18 proves the passage through the compact central factor and its
finite covering to $U(p,q)$. Thus no compact-center or measure
normalization hypothesis is suppressed here.

For every literal representation at this parameter and every nonzero
vector, its formal degree is this same number. The lower bound
$(1-\delta)d_n(\lambda)\le d_n(\lambda)$ is immediate. For the upper
bound use (DSZ.2), $3pq\ge1$, and the lowest-type dimension just
identified. This gives exactly
\[
 (1-\delta)d_n(\lambda)\le d(\sigma)
 \le4\dim(\sigma_{\rm lowK})(3pqA)^{pq}.                       \tag{DSZ.4}
\]
\emph{Source: Z04--Z06.}
% END DSZ2

% BEGIN DSZ3
\paragraph{DSZ3: identify the stored degree and close the actual $h_3$ slot.}
\label{claim:DSZ3}
For the original $D$, the row has length
$r=2r_S+k+1=(k-1)+m$, signature
$(p,q)=(k-1+m/2,m/2)$, sign vector
$\texttt{rowSgn}$ and parameter
\[
 \lambda_\rho=(v_0,\ldots,v_{k-2},-b_0,\ldots,-b_{m-1}).
                                                               \tag{DSZ.5}
\]
The original finite reindexing identifies these with
$\texttt{epsSH}$ and $\texttt{lamSH}$; it changes neither a root
difference nor a product. The \texttt{HCParam} clauses give precisely
regular integrality. Also $p,q\ge1$ because $m\ge2$ is even.

Assume the two original ties \texttt{cdScaleTie} and
\texttt{dsigTie}. If (DSZ.5) has a noncompact gap satisfying
(DSZ.3), DSZ2 constructs a representation and proves its
\texttt{IsFdeg} identity for the same scaled measure. Apply
\texttt{dsigTie} to that representation and that degree. It yields
the exact, stronger equality
\[
                    D.\mathrm{dsig}(\rho)=d_r(\lambda_\rho).
                                                               \tag{DSZ.6}
\]
This step is essential: a positive stored number or a degree of an
unrelated representation would not give (DSZ.6).

At the source threshold \texttt{spectralGapLowerBound}, (DSZ.3)
is exactly the displayed hypothesis. Substituting (DSZ.6) and
DSZ1 into (DSZ.4) proves the original
\texttt{SigmaDegreeTwoSidedBoundOfGap}, including its quantifiers
over every $\delta$ and every place. In a \texttt{FamP}, both ties
are fields of $\mathcal F.\texttt{ties}(x,i)$ for every $x,i$.
Hence this predicate holds at every member, with eventual threshold
$0$. Its \texttt{Closed} version on any class has
$\epsilon_1=1$ and the zero transported bound $K$. These bounds
are chosen before a family. No new existential row selector is
required for this slot. This closes \texttt{TASK\_GL3\_GAP\_DEGREE}
relative to the already recorded HC-Deg application.
\emph{Source: Z06--Z09.}
% END DSZ3

% BEGIN DSZ4
\paragraph{DSZ4: read the needed inequalities from the actual family.}
\label{claim:DSZ4}
Fix a supplied family, so $0<\epsilon<1/52$, and put
$C_0=\max\{5,C_{YL}+1\}$,
$\Delta=D.s.\mathrm{Del}$,
$S=4k\Delta=\mathcal L.\mathrm{Sc}$,
$\eta=\mathcal L.\mathrm{etaSz}$, and
$w_{j\rho}=\texttt{wj}(D,j+1,\rho)$.
Here $\mathcal L=\mathcal F.L(x,i)$; this notation distinguishes
that record from the scalar $L=x^\epsilon$.
The following checks use its actual \texttt{params},
\texttt{sizes}, and \texttt{hcW} identities. They do not replace its
coordinates by a new equally spaced model.

The rank windows give, uniformly in $i$ and eventually in $x$,
\[
 x^{1/4}\le m\le x,\quad f,k\le x^5,\quad k\ge3x,
 \quad m\le k,\quad m^2L\le k,\quad L\ge1.                    \tag{DSZ.7}
\]
For example the lower exponent for $m$ is
$1/2-\epsilon/16-\epsilon/100>1/4$; its upper exponent is
$1/2-\epsilon/16+\epsilon/100<1$. The actual upper exponent for
$k$ is $9/2+10\epsilon+\epsilon/100<5$, and its lower exponent
exceeds $2$. The $n$ bound and $f=2n$ give $f\le x^5$ for
$x\ge3$. Increasing the threshold so that $5\log x\le L$ also
gives $k\le B$. The inequality $m^2L\le k$ is a stated size clause.
These comparisons retain the $\epsilon/100$ margins.

The host discriminant lower bound, as in QW8, gives $v>\pi$.
For the upper bound, put $Q=Q_Sqd\ge1$. The same host bounds
give $v\le Q$, while \texttt{SizeCond} says
$Q^{36}t^{150}\le\sqrt B$ with $t\ge1$.
Thus $v\le B$ and $\log v\le L$.
The auxiliary-prime parameters give one $Y$ with
$(kB)^8\le Y\le g_i\le2Y$ and $\log Y\le C_{YL}L$.
Therefore
\[
                  8L\le\log g_i\le C_0L.                    \tag{DSZ.8}
\]
The size bound $N(m)m^{10}\le\Delta$ gives $\Delta\ge1$.

The exact cloud identities give, for the indicated indices,
\[
 0\le v_i\le4k\Delta,\quad v_i-v_{i+1}=4\Delta,
 \qquad u_l=v_{l-1}-\tfrac12-\xi_l\Delta,
 \quad \xi_l\in\{0,1\},\quad1\le l\le k-2.                 \tag{DSZ.9}
\]
At a non-top place the actual endpoint clauses give
$4k\Delta\le u_0-v_0\le8k\Delta$ and
$u_{k-1}=v_{k-2}-1/2$. At a top place the literal endpoint formulas
are $u_0=a_T+(k-1)/2$ and $u_{k-1}=-b_T-(k-1)/2$.
The window $a\le a_T\le a+a/(kx^{26/2})$ gives
$a_T\le4a/3$. The bound $b_T\le5B'/4$ established below,
together with $k\le B$, $B'\le B$ and $3B\le a$, therefore gives
$u_0\le4a/3+B/2\le2a$ and
$u_{k-1}\ge-5B/4-B/2\ge-a$.

Here are the bounds which justify those window estimates and the
remaining weight clauses. Put
$X'=\exp(x^{\epsilon/3})$, $B'=xB/X'$ and $w''=B'/(xn)$.
For fixed $\epsilon>0$, eventually $X'\ge2x$; hence
$0<w''\le B'\le B/2$.
The actual size inequalities include
\[
 \eta\le B/(t^3fk^3)\le B/27,\qquad
 t\eta\le w''/x,\qquad kS\le w'',\qquad t\ge3.
                                                               \tag{DSZ.10}
\]
Consequently $(x+1)\eta\le w''$ and $xS\le w''/3$.
The top window has $B'\le b_T\le5B'/4$.
For its upper bound, the source's lower-window width is at most
$B'/4$ by \texttt{SizeRho}; all quantities discarded from that
size inequality are nonnegative. At top, thick and thin places
the respective upper bounds on $w_j$ are
$b_T+B'/x+\eta$, $B+\eta+1$, and $w''+\eta$.
Their lower bounds are respectively $B'$, $B$, and $w''$.
Since $B\ge e$, the additional $1$ in the thick window is also
absorbed by $2B$. Thus in every case
\[
 w''\le w_j\le2B,\quad b_j<w_j,\quad0<w_j-b_j\le\eta,
 \quad |w_j-b_j|\le b_j/x.                                    \tag{DSZ.11}
\]
Indeed $b_j=w_j-(w_j-b_j)\ge w''-\eta\ge x\eta$.
Also $xS+\eta\le w''$ gives $4k\Delta\le b_0/x$.
The original $e_j,b_{j+1}$ separation gives $w_j<b_{j+1}$ and
$b_{j+1}-b_j\ge1$. Original interlacing gives
$u_i>v_i>u_{i+1}$ and $u_{k-1}>-b_0$.
Every $u_i-u_j$ with $i<j$ is at least one, by the original
integral column weight, not just by strict positivity.

Finally \texttt{nodes\_gap} and \texttt{SizeNodes} give
\[
 b_{j+1}-b_j\ge\eta/m
       \ge x^{\epsilon/100}r^{19}f^5L
       \ge r^{19}f^5L/C_0.                                   \tag{DSZ.12}
\]
Equations (DSZ.7)--(DSZ.12), with the displayed cloud and endpoint
identities, check every clause of the actual \texttt{DataB} needed
below. In particular they check QW4--QW6's \texttt{PlaceData} on
every word and place; $w_j\ge1$ follows from $b_0\ge4k\Delta x$.
All choices made were thresholds for this given family, uniform in
$i$, words and places. No bound on a newly chosen $\Delta$ is assumed.
\emph{Source: Z10--Z15.}
% END DSZ4

% BEGIN DSZ5
\paragraph{DSZ5: the exact formal-degree polynomial on every late member.}
\label{claim:DSZ5}
In (DSZ.5) all cloud coordinates have positive sign. The signs of
the negative bottom coordinates alternate. Two opposite-sign bottom
coordinates are separated by at least one adjacent difference in
(DSZ.12). A cloud coordinate and a negative-sign bottom coordinate
are separated by at least the first bottom difference: the first
negative-sign bottom index is $1$, $v_i\ge0$, and $b_0>0$.
It follows that the actual row's noncompact gap is at least
$r^{19}f^5L/C_0$.

The threshold (DSZ.3) for $\delta=1/3$ is at most
$9\cdot8192\,r^9$: use $p,q,\min(p,q)\le r$ and
$\log r\le r$ (this estimate includes harmless slack).
Once $L\ge9\cdot8192 C_0$, the gap just proved exceeds this
threshold, since $r\ge1$ and $f\ge1$.
Also $r\ge2$ gives $e^{-r}\le1/3$.
DSZ3 therefore gives (DSZ.6) simultaneously for every original
place and index of every sufficiently late member of the family.
In particular the definition of $D_*$ becomes exactly
\[
 D_* =\mu_r I_{\rm row}\prod_\rho d_r(v_\rho,-b_\rho)
                                                               \tag{DSZ.13}
\]
whenever its original volume satisfies $\operatorname{vol}(Y_0)=\mu_r$.
The corresponding $d_*(u)=\mu_k I_{\rm col}\prod_\rho d_k(u_\rho)$
is already the original definition. The \texttt{FactDstar} and
\texttt{FactDuStar} family fields identify these with the stored
positive $D_*,d_*(u)$ in the actual prime window.
\emph{Source: Z09, Z14--Z17.}
% END DSZ5

% BEGIN DSZ6
\paragraph{DSZ6: the original mass, index and factorial identity.}
\label{claim:DSZ6}
Assume the actual volume identity in DSZ5 and the mass-ratio theorem
already reconstructed in MF19 and RSL1--RSL3, with their named
classical inputs retained. Put
\[
 X_u=\frac1{fk}\log\frac{D_*}{d_*(u)},\quad
 \ell=\log(v/\pi),\quad G=\log(g_0g_1),\quad
 A_m=(m-1)k+\frac{m(m-1)}2,\quad
 B_m=\frac{m(2k+m-1)}2.
\]
By QW2 the exact index ratio, before weakening any Euler factors, is
\[
 \frac{I_{\rm row}}{I_{\rm col}}
  =(g_0g_1)^{fB_m}
       \prod_{i=0}^1\prod_{j=k+1}^{r}(1-g_i^{-j})^f.
\]
The equality follows by dividing the two exact finite-group
cardinalities in equations (QW.4)--(QW.6), not just their inequalities.
QW3 telescopes the actual positive masses, and its factorial
factors cancel those in (DSZ.13) and $d_*(u)$ exactly as in
(QW.15). Thus
\begin{align*}
 fkX_u={}&fA_m\ell-\frac{f(m-1)}2\log(2v)
     +\sum_{M=k+1}^{r}\log L_E(M)+fB_mG\\
 &+f\sum_{i=0}^1\sum_{j=k+1}^{r}\log(1-g_i^{-j})
     +\sum_\rho\log\frac{V_r(v_\rho,-b_\rho)}{V_k(u_\rho)}.
                                                               \tag{DSZ.14}
\end{align*}
Every logarithm is of a positive number; this was proved for the
Euler product in QW1, the masses in IG9/MF, the indices in QW2,
and the Vandermondes by the actual interlacing.
This is the source's \texttt{LemmaFiveOneIdentity} after the
stronger degree equality (DSZ.6). The total factorial contribution
is zero, rather than an error estimated at size $k^2$.
\emph{Source: Z16--Z19.}
% END DSZ6

% BEGIN DSZ7
\paragraph{DSZ7: a short upper Vandermonde comparison on the actual cloud.}
\label{claim:DSZ7}
Fix one actual word and place satisfying DSZ4. Let $n=k-2$ and
compare the interior column coordinates $u_1,\ldots,u_{k-2}$
with $v_0,\ldots,v_{k-3}$. From (DSZ.9), for $1\le i<j\le k-2$,
\[
 u_i-u_j\ge\Delta(4(j-i)-1),\qquad
 v_{i-1}-v_{j-1}=4\Delta(j-i).
\]
Therefore, with empty products interpreted as one,
\[
 \log\frac{V_n(v_0,\ldots,v_{k-3})}
               {V_n(u_1,\ldots,u_{k-2})}
 \le\sum_{d=1}^{n-1}(n-d)\log\frac{4d}{4d-1}
 \le\frac{k}{3}(1+\log k).                                  \tag{DSZ.15}
\]
For the last inequality, use
$\log(1+t)\le t$, $4d-1\ge3d$, and
$\sum_{d=1}^{n-1}d^{-1}\le1+\log k$.
This harmonic sum is the reason the comparison costs $k\log k$
rather than $k^2$.

The remaining $k-2$ factors in $V_{k-1}(v)$ involve its last
coordinate and are at most $4k\Delta\le2B$.
The $(k-1)m$ cloud-to-bottom factors are at most $4B$, since
$0\le v_i\le4k\Delta\le2B$ and $0<b_j<2B$.
The $m(m-1)/2$ bottom factors are at most $2B$.
Every omitted endpoint factor in $V_k(u)$ is at least one by
DSZ4. Combining these facts gives
\begin{align*}
 \log\frac{V_r(v,-b)}{V_k(u)}
 \le{}&\frac{k}{3}(1+\log k)+(k-2)\log(2B)\\
     &+(k-1)m\log(4B)+\frac{m(m-1)}2\log(2B).
                                                               \tag{DSZ.16}
\end{align*}
Eventually $L\ge3$ and $\log k\le5\log x\le L$.
Then $\log(2B),\log(4B)\le2L$; using $2\le m\le k$,
the four terms in (DSZ.16) are bounded respectively by
$kmL/3$, $kmL$, $2kmL$, and $kmL$. In particular
\[
                 \log\frac{V_r(v,-b)}{V_k(u)}\le5kmL.
                                                               \tag{DSZ.17}
\]
This is an independent proof for the literal coordinates. It is
stronger than needed for the size bound and replaces only that
use of the source's more elaborate one-place upper estimate.
\emph{Source: Z12, Z14, Z20--Z21.}
% END DSZ7

% BEGIN DSZ8
\paragraph{DSZ8: two-sided logarithmic bounds with one family threshold.}
\label{claim:DSZ8}
Under the hypotheses of DSZ6, apply (DSZ.14).
We have $0\le\ell\le L$, $\log(2v)>0$, $G\le2C_0L$,
and $\sum_{M=k+1}^r\log L_E(M)\le fm$ by QW1.
The local-factor logarithms are nonpositive.
Also $A_m\le3km/2$, $B_m\le3km/2$, and (DSZ.17) bounds
the final place sum by $5fkmL$. Since $k\ge5$ and $L\ge1$,
$fm\le fkmL/2$. Dividing by $fk>0$ gives
\[
                         X_u\le(7+3C_0)mL.                   \tag{DSZ.18}
\]

For the lower bound use QW6 with the same actual coordinates,
$C_0$, and the exact (DSZ.13), which implies its weaker
$(2/3)^f$ hypothesis. All its remaining hypotheses were checked
in DSZ4. Put
$C_\epsilon=30+(2+2C_0)/\epsilon$.
The average of $\log(w_jv/\pi)$ is nonnegative, and (DSZ.8)
gives $mG\ge16mL$. Since $m^2L\le k$ and $L\ge1$,
\[
 X_u\ge16mL-C_\epsilon(L+m^2L/k)
       \ge16mL-2C_\epsilon L.
\]
Choose the family threshold also so that $m\ge2C_\epsilon$,
which follows from $m\ge x^{1/4}$. We obtain, simultaneously
for every original index and word,
\[
                    15mL\le X_u\le(7+3C_0)mL.                \tag{DSZ.19}
\]
The family-dependent constants $C_0,C_\epsilon$ precede this
eventual threshold, as the definition of \texttt{TuSize} permits.
The original volume identity may itself have a family threshold;
take its maximum with the finitely many thresholds above.
The argument does not require a uniform value of this final
threshold for all possible families.
\emph{Source: Z17--Z22; QW1--QW6.}
% END DSZ8

% BEGIN DSZ9
\paragraph{DSZ9: obtain the literal prime-size predicate with the original constant.}
\label{claim:DSZ9}
Let $\Gamma_u$ be the actual $\texttt{GamP}$.
The original \texttt{Gam\_eq} and \texttt{SizeKRel} inequality
give $\log\Gamma_u\le fk$; also $\Gamma_u>1$ is an original
field. Thus the actual prime-window center satisfies
\[
 T_u^*=(D_*/d_*(u))^{1/(fk)}\Gamma_u^{-2/(fk)},\qquad
 X_u-2\le\log T_u^*\le X_u.                                  \tag{DSZ.20}
\]
The source selects its original prime in
$T_u^*\le T_u\le(1+1/(fk))T_u^*$.
Since $fk\ge1$, taking logarithms and using
$\log(1+1/(fk))\le1$ gives
\[
 X_u-2\le\log T_u\le X_u+1.
\]
By (DSZ.19), $m\ge2$ and $L\ge1$, the lower bound is at least
$mL/2$. The upper bound is at most $(8+3C_0)mL$.
Consequently
\[
 e^{mL/2}=B^{m/2}\le T_u,
 \qquad \log T_u\le(8+3C_0)mL.                               \tag{DSZ.21}
\]
These are the precise lower predicate \texttt{TuGe} and upper
inequality in \texttt{TuSize}. In particular they prove the
source conclusion whenever
\[
 c_T(C_{YL})=\max\{11,16+3\max(5,C_{YL}+1)\}\le C_{Tu},       \tag{DSZ.22}
\]
using its original, larger constant without changing any chosen
prime or any family record. Only a threshold in $x$ was enlarged.

We have proved: every original supplied family with eventual
$\operatorname{vol}(Y_0)=\mu_r$ satisfies
$c_T(C_{YL})\le C_{Tu}\Longrightarrow\texttt{TuSize}$.
All other premises were original family fields or the already
matched mass and formal-degree inputs. In particular no unresolved
Vandermonde variation, rank formula, local Hecke estimate, or
small-level rank theorem is used in this proof.
\emph{Source: Z15, Z17, Z23--Z25.}
% END DSZ9

% BEGIN DSZ10
\paragraph{DSZ10: the design-class row supplier and its quantifiers.}
\label{claim:DSZ10}
Take exactly the original construction class
$\mathcal A=\texttt{cls}(Q,\mathrm{Lk},E)$ and its common selector.
A \texttt{Built} witness is a context $K$, a point $\beta$, a
linked bundle $b'$ and an equality of the entire dependent triple
$(D,O,b)=(\beta.D,\beta.O,b')$.
At $\beta$, the stored core carries the actual
\texttt{RealTied} witness. The subsequent \texttt{setDiv} operation
changes only the division lattices and their prescribed heights;
it preserves the host, Hermitian spaces, degree function, compact-dual
measure and volume. Its four \texttt{RealTied} clauses therefore
remain the same. Transport along the displayed dependent equality
gives \texttt{RealTied} for the original $D$ of the selected bundle.
This is the content of \texttt{built\_val}; it is a projection and
transport of the supplied construction witness, not a new
existence assertion about a core.

Every member of $\mathcal A$ supplies such a \texttt{Built}
witness at every $x>\mathcal F.x_1$ and every actual $i$, by its
literal \texttt{QOf} and \texttt{Link} clauses. The parameter
positivity guard is discharged by that same Link's positivity
field. RSL1--RSL3 apply to this \texttt{RealTied} and give
$\operatorname{vol}(Y_0)=\mu_r$ for those very objects, with
the already recorded GV Tamagawa and MF Hasse--Minkowski
applications retained. Thus DSZ9 applies on this class and on
every subclass obtained by adjoining \texttt{WithChN},
\texttt{WithSm}, \texttt{WithTw}, or the target point condition.

For the source's \texttt{Closed} assertion choose
$\epsilon_1=1$ and the zero transported bound $K$, before the
family and $CH$ are given. Once a family in the specified class
is supplied, forget the added class conditions, project the
\texttt{Built} witness, and apply DSZ9 with its permitted
family-dependent eventual threshold. The \texttt{Large}
premise is unused. This proves the exact
\[
 \texttt{Closed}\bigl(\mathcal C_{\rm des},\,
       c_T(C_{YL})\le C_{Tu}\Longrightarrow\texttt{TuSize}\bigr)
                                                               \tag{DSZ.23}
\]
for the one construction class and the original function
$c_T(\epsilon,y)=c_T(y)$ in (DSZ.22).
The same proof also works for an arbitrary class $\mathcal A$ and
selector $\sigma$ supplied with
$\texttt{IsSel}(Q,\mathrm{Lk},\mathrm{True},\mathcal A,\sigma)$:
the definition of \texttt{IsSel} supplies the required \texttt{Link}
and the \texttt{Built} witness for each literal selected bundle.
The proof above then applies without using the special definition
of $\texttt{cls}$. This establishes the universal premise of
\texttt{Row.of\_sel}, and hence the original
\texttt{row\_hTu\_cTmin} row with the trivial extra bundle predicate
$E=\mathrm{True}$.
The finite conjunction with other rows therefore imposes no new
bundle duty here. DEX6--DEX8 allow this proved row to be joined
before one family is selected.

\emph{Audit consequence.} The former \texttt{DESIGN\_TUSIZE\_CLOSED}
leaf is discharged on the selected construction-class route.
The independent compatible \texttt{DesignChoices} producer remains;
(DSZ.23) does not construct a member of that class. The result also
does not assert a volume identity for arbitrary \texttt{FamP}
records outside the construction class. These restrictions are
part of the theorem just proved, not new assumptions to be
silently supplied by positivity.
\emph{Source: Z26--Z29.}
% END DSZ10
